\chapter{Flags, convex bodies and algebraic directions}
\label{chap:flag-bodies}

A flag converts multiplication of sections into addition of integer vectors. Gauduchon degrees bound these vectors linearly in the tensor power, so their normalized convex hull is compact. Counting its lattice points gives the growth of the section algebra, with the dimension and degree period determined by its field of section ratios and its nonzero degrees. A point blowup supplies the required flag on every compact complex manifold while preserving its sections.

Full-dimensional growth has a geometric consequence: independent valuation directions give algebraically independent meromorphic functions. The resulting projective modification bounds Hermitian Monge--Amp\`ere masses. On a manifold where bounded mass fails, the section bodies therefore have smaller dimension; the Hopf example computes this dimension and its growth coefficient explicitly.

\section{A complete compact flag}\label{sec:compact-flag}

A complete flag records successive vanishing orders along divisors. In local coordinates these are the first exponent of a Taylor series; the leading coefficient then identifies the number of distinct values with the dimension of the section space.

Choose $p\in X$ and fix
\[
 X=Y_0\supset Y_1\supset\cdots\supset Y_n=\{p\},
\]
where every $Y_i$ is smooth, connected and compact, and is a prime divisor in $Y_{i-1}$. Let $L$ be a holomorphic line bundle. For a nonnegative integer $m$ and a nonzero section $s\in H^0(X,L^m)$, define $\nu_1(s)$ to be its vanishing order along $Y_1$. Divide by the corresponding power of the canonical section of $\mathcal O_X(Y_1)$ and restrict to $Y_1$. The result is a nonzero section of $L^m(-\nu_1Y_1)|_{Y_1}$. Repeat the procedure down the flag. The resulting vector
\[
 \nu(s)=(\nu_1(s),\ldots,\nu_n(s))\in\mathbb N^n
\]
is its flag valuation.

\Needspace{9\baselineskip}
Near $p$, choose coordinates adapted to the flag and a nonvanishing frame of $L$. This construction is the lexicographically first exponent of the Taylor expansion of the local coefficient of $s$. In particular, for nonzero sections of the indicated degrees,
\[
 \nu(st)=\nu(s)+\nu(t),
 \qquad \nu(s+t)\geq\min\{\nu(s),\nu(t)\}
\]
when $s,t$ have the same degree and $s+t\neq0$. If their valuations are different, equality holds in the second formula. The inequality uses lexicographic order.

\begin{lemma}\label{lem:flag-one-dimensional-leaves}
For every finite-dimensional vector subspace $V\subset H^0(X,L^m)$,
\[
 \#\{\nu(s):0\neq s\in V\}=\dim V.
\]
\end{lemma}
\begin{proof}
Sections with distinct valuations are linearly independent: in a nontrivial finite linear combination, the least valuation occurs in exactly one summand and cannot cancel. Thus there are at most $\dim V$ values.

For the reverse bound, proceed by induction on $\dim V$. The nonempty set of values has a least element in the lexicographic order on $\mathbb N^n$. The coefficient of that monomial is a nonzero linear functional on $V$. Its kernel has codimension one, and every nonzero member of the kernel has strictly larger valuation. By induction the kernel contributes $\dim V-1$ distinct values, in addition to the least one. The zero-dimensional case is immediate.
\end{proof}

\section{Gauduchon bounds for the valuation}

The residual section at each stage of the flag has an effective zero divisor. Pairing it with a Gauduchon metric bounds the next vanishing order in terms of the degree and the preceding orders. This recursion gives a bound uniform over all sections and tensor powers.

\begin{theorem}\label{thm:gauduchon-linear-flag-bound}
There are constants $C_1,\ldots,C_n\geq0$, depending only on $L$ and the chosen flag, such that
\[
 0\leq\nu_i(s)\leq C_i m
 \quad(1\leq i\leq n)
\]
for every $m\geq1$ and every nonzero $s\in H^0(X,L^m)$.
\end{theorem}
\begin{proof}
For each $i$, equip $Y_{i-1}$ with a Gauduchon metric $\gamma_i$; when it is a curve, any Hermitian metric will do. Put $d_i=n-i+1$ and
\[
 a_i=\int_{Y_i}\gamma_i^{d_i-1}>0,
 \qquad A_i=\deg_{\gamma_i}(L|_{Y_{i-1}}).
\]
For $i=n$ the integral defining $a_i$ is the integral of $1$ over the point $p$, so $a_n=1$.

\Needspace{6\baselineskip}
For $j<i$, also put
\[
 B_{ij}=\deg_{\gamma_i}
 \left(\mathcal O_{Y_{j-1}}(Y_j)|_{Y_{i-1}}\right).
\]
All these are fixed real numbers.

\Needspace{8\baselineskip}
At stage $i$ of the valuation algorithm, the residual section belongs to the line bundle
\[
 M_{i-1}=L^m|_{Y_{i-1}}\otimes
     \bigotimes_{j<i}\mathcal O_{Y_{j-1}}(-\nu_j(s)Y_j)|_{Y_{i-1}}.
\]
\Needspace{6\baselineskip}
Its zero divisor contains $\nu_i(s)Y_i$. Poincar\'e--Lelong, tested against $\gamma_i^{d_i-1}$, therefore gives
\begin{equation}\label{eq:flag-degree-recursion}
 a_i\nu_i(s)\leq mA_i-\sum_{j<i}B_{ij}\nu_j(s).
\end{equation}
For $i=1$ this bounds $\nu_1(s)$ by $m|A_1|/a_1$. If the bounds are known for the preceding indices, the right side is at most
\[
 m\left(|A_i|+\sum_{j<i}|B_{ij}|C_j\right).
\]
Divide by $a_i$ to choose $C_i$ and complete the induction.
\end{proof}

\section{The body of a graded section algebra}\label{sec:graded-flag-body}

The additive valuation vectors of a graded algebra lie in a linearly bounded cone. Its height-one slice is the convex body whose lattice points count sections.

Let $W_\bullet=\bigoplus_{m\geq0}W_m$ be a graded subalgebra of the section ring of $L$, with $W_0=\mathbb C$. One example is the bounded weighted section algebra $\mathcal W_\bullet(u)$ of \cref{chap:sections}.

Set
\[
 S=\{(0,0)\}\cup
       \{(\nu(s),m):m\geq1,\ 0\neq s\in W_m\}.
\]
It is an additive semigroup. Indeed the product of two nonzero sections is nonzero on the connected manifold, and valuations add. Choose constants $C_i$ from \cref{thm:gauduchon-linear-flag-bound}. Then $S$ lies in the closed cone
\[
 \{(x,t)\in\mathbb R^n\times\mathbb R:t\geq0,\ 0\leq x_i\leq C_i t\}.
\]
The only point of this cone at height zero is the origin. Define
\[
 \Delta(W_\bullet)=
 \overline{\operatorname{conv}}
 \{m^{-1}\nu(s):m\geq1,\ 0\neq s\in W_m\}.
\]
This is a compact convex subset of $\mathbb R^n$ when the set is nonempty. If every positive-degree piece is zero, set $\Delta(W_\bullet)=\varnothing$. Equivalently it is the height-one slice of the closed cone generated by $S$. Write $\operatorname{vol}_{\mathbb R^n}$ for Euclidean volume.

\Needspace{10\baselineskip}
\begin{proposition}\label{prop:flag-body-upper-count}
For every such algebra,
\[
 \limsup_{m\to\infty}m^{-n}\dim W_m
 \leq\operatorname{vol}_{\mathbb R^n}\Delta(W_\bullet).
\]
In particular a body with empty interior gives $\dim W_m=o(m^n)$.
\end{proposition}
\begin{proof}
Let $S_m=\{\nu(s):0\neq s\in W_m\}$. Its cardinality is $\dim W_m$ by \cref{lem:flag-one-dimensional-leaves}, and $m^{-1}S_m\subseteq\Delta(W_\bullet)$. Attach the half-open cube $m^{-1}(a+[0,1)^n)$ to each $a\in S_m$. The cubes are disjoint, their total volume is $m^{-n}\dim W_m$, and their union is contained in
\[
 \Delta(W_\bullet)+[0,m^{-1}]^n.
\]
These compact neighborhoods decrease to $\Delta(W_\bullet)$, so their volumes converge to its volume. This proves the inequality. A compact convex set has positive $n$-dimensional volume exactly when it has nonempty interior, proving the last statement. The empty-body case is immediate.
\end{proof}

The semigroup growth theorem recalled in \cite[appendix ``Almost semigroups'', Kaveh--Khovanskii theorem]{SGPT} gives an exact count when the generated lattice is full. It states that if an additive semigroup in $\mathbb Z^{n+1}$ lies in a closed convex cone in $\mathbb R^n\times\mathbb R_{\geq0}$ with compact height-one slice and with only the origin at height zero, and generates $\mathbb Z^{n+1}$ as a group, then its degree-$m$ cardinalities divided by $m^n$ converge to the Euclidean volume of the height-one slice of its own closed cone.

\begin{corollary}\label{cor:flag-body-primitive-count}
If the subgroup generated by $S$ is all of $\mathbb Z^{n+1}$, then
\[
 \lim_{m\to\infty}m^{-n}\dim W_m
 =\operatorname{vol}_{\mathbb R^n}\Delta(W_\bullet).
\]
\end{corollary}
\begin{proof}
The cone bound and its height-zero condition were proved above, and the subgroup hypothesis is precisely the one in the stated semigroup theorem. Apply it and use \cref{lem:flag-one-dimensional-leaves}.
\end{proof}

\section{Moishezon manifolds and bounded mass}

Independent directions in the body give meromorphic functions with independent valuations. Their algebraic independence produces a projective modification, on which a fixed K\"ahler volume bounds the masses of all Hermitian potentials pulled back from $X$.

\begin{proposition}\label{prop:full-flag-body-meromorphic}
If $\Delta(W_\bullet)$ has nonempty interior in $\mathbb R^n$, then $X$ has $n$ algebraically independent meromorphic functions.
\end{proposition}
\begin{proof}
The normalized valuation vectors of actual sections affinely span $\mathbb R^n$; otherwise their closed convex hull would lie in a proper affine subspace. Choose sections $s_0,\ldots,s_n$ of positive degrees $m_0,\ldots,m_n$ with affinely independent normalized vectors. Let $M$ be a common multiple of the degrees and put $t_i=s_i^{M/m_i}\in W_M$. The quotients
\[
 f_i=t_i/t_0\quad(1\leq i\leq n)
\]
are global meromorphic functions. The valuation extends to such quotients by subtraction, and
\[
 \nu(f_i)=M\left(\frac{\nu(s_i)}{m_i}
                       -\frac{\nu(s_0)}{m_0}\right).
\]
These $n$ vectors are linearly independent. Distinct monomials in the $f_i$ consequently have distinct valuation vectors. In any nonzero polynomial expression in the $f_i$, the term of least lexicographic valuation cannot cancel. Thus no nonzero polynomial vanishes on $(f_1,\ldots,f_n)$, proving algebraic independence.
\end{proof}

We use the classical projective-modification theorem: a compact irreducible complex space of dimension $n$ with $n$ algebraically independent meromorphic functions admits a modification from a smooth projective manifold. This form of Moishezon's theorem is recorded in \cref{sec:history-sections}.

\begin{lemma}\label{lem:kahler-modification-bounded-mass}
If a compact complex manifold $X$ admits a modification $\pi:Z\to X$ from a compact K\"ahler manifold, then $X$ has the bounded mass property.
\end{lemma}
\begin{proof}
Fix a Hermitian form $\omega$ on $X$ and a K\"ahler form $\kappa$ on $Z$. Choose $C>0$ with $\pi^*\omega\leq C\kappa$. By \cref{prop:bmp-metric}, it suffices to bound the masses of smooth $u\in\PSH(X,\omega)$. Their pullbacks satisfy
\[
 0\leq\pi^*\omega+\ddc(u\circ\pi)
 \leq C\kappa+\ddc(u\circ\pi).
\]
Taking top powers and integrating gives
\[
 \begin{split}
 \int_X(\omega+\ddc u)^n
 &=\int_Z(\pi^*\omega+\ddc(u\circ\pi))^n\\
 &\leq\int_Z(C\kappa+\ddc(u\circ\pi))^n
 =C^n\int_Z\kappa^n.
 \end{split}
\]
The first equality follows by change of variables on the biholomorphic locus: smooth top-degree forms give the exceptional analytic sets zero mass. The inequality follows by expanding in the smooth semipositive difference $C\kappa-\pi^*\omega$. Since $\kappa$ is closed, Stokes' theorem gives the last equality. The bound is independent of $u$, proving bounded mass.
\end{proof}

\begin{theorem}\label{thm:bmp-failure-flag-degeneracy}
If $X$ fails the bounded mass property, then for every complete smooth compact flag, every holomorphic line bundle $L$ and every graded section subalgebra $W_\bullet$ as above,
\[
 \operatorname{vol}_{\mathbb R^n}\Delta(W_\bullet)=0,
 \qquad \dim W_m=o(m^n).
\]
\end{theorem}
\begin{proof}
A body of nonempty interior would give $n$ algebraically independent meromorphic functions by \cref{prop:full-flag-body-meromorphic}. The stated Moishezon theorem would give a smooth projective, hence K\"ahler, modification of $X$. This contradicts \cref{lem:kahler-modification-bounded-mass}. The body therefore has empty interior, and \cref{prop:flag-body-upper-count} gives both conclusions.
\end{proof}

\section{Intrinsic dimension, lattice and growth}

A body of zero volume in $\mathbb R^n$ can still detect growth of lower order. The group generated by the valuation semigroup determines a degree period and a lattice in which the body is full-dimensional, and counting in this lattice gives the growth exponent and the leading coefficient. Ratios of sections identify the exponent with a transcendence degree, and since the coefficient is a limit of graded dimensions, it does not depend on the flag.

Throughout this section, some positive-degree piece of $W_\bullet$ is nonzero. Let $G\subset\mathbb Z^{n+1}$ be the subgroup generated by $S$, and let $\operatorname{pr}_{n+1}$ be the projection to the last coordinate. Define the \emph{degree period} $e\geq1$ by
\[
 \operatorname{pr}_{n+1}(G)=e\mathbb Z.
\]
Thus $e$ is the greatest common divisor of the degrees $m$ with $W_m\neq0$. Set
\[
 \Lambda=G\cap(\mathbb R^n\times\{0\}),
 \qquad d=\operatorname{rank}\Lambda.
\]
We regard $\Lambda$ as a lattice in its real span, a $d$-dimensional subspace of $\mathbb R^n\times\{0\}$. Its lattice measure assigns volume one to a fundamental parallelepiped. On each translate of this span we use the translated measure, denoted by $\operatorname{vol}_\Lambda$; when $d=0$, a point has volume one.

\begin{theorem}\label{thm:intrinsic-flag-count}\index{Convex body!intrinsic lattice volume}\index{Degree period}
The affine dimension of $\Delta(W_\bullet)$ is $d$. Its lattice volume is finite and strictly positive, and
\[
 \lim_{q\to\infty}\frac{\dim W_{eq}}{(eq)^d}
 =\operatorname{vol}_\Lambda\Delta(W_\bullet).
\]
If $e$ does not divide $m$, then $W_m=0$. In particular, $W_m\neq0$ for all sufficiently large multiples $m$ of $e$.
\end{theorem}
\begin{proof}
Choose $g_0=(b,e)\in G$ and an integral basis $\lambda_1,\ldots,\lambda_d$ of $\Lambda$. Every element of $G$ has a unique expression
\[
 \sum_{i=1}^d a_i\lambda_i+qg_0,
 \qquad (a,q)\in\mathbb Z^d\times\mathbb Z.
\]
These coordinates identify $G$ with $\mathbb Z^{d+1}$. Let $S'$ be the image of $S$ under this identification. It generates $\mathbb Z^{d+1}$, and the height of the image of $(x,m)\in S$ is $q=m/e\geq0$.

The points $(m^{-1}x,1)$, for $(x,m)\in S$ with $m>0$, span the $(d+1)$-dimensional real vector space $\operatorname{span}_{\mathbb R}G$. Since they all have height one, their affine span has dimension $d$. Taking the convex hull and the closure does not change this affine span, which proves the dimension assertion.

Intersect the cone containing $S$ with $\operatorname{span}_{\mathbb R}G$. The linear change of coordinates carries this intersection to a closed cone with compact height-one slice whose only point at height zero is the origin. Indeed, a vector $(x,m)$ with $|x|\leq Cm$ has horizontal coordinates bounded by a constant times $m$, hence by a constant times $q$. The height-one body of $S'$ is
\[
 \Delta'=e\bigl(\Delta(W_\bullet)-b/e\bigr),
\]
where the right-hand side is expressed in the basis $\lambda_1,\ldots,\lambda_d$. It is full-dimensional in $\mathbb R^d$, and therefore
\[
 \operatorname{vol}_{\mathbb R^d}\Delta'
 =e^d\operatorname{vol}_\Lambda\Delta(W_\bullet)>0.
\]
Apply the semigroup counting theorem of \cref{sec:graded-flag-body} to $S'$, which generates its full lattice. The points of $S'$ at height $q$ correspond bijectively to the valuations of the nonzero elements of $W_{eq}$, and by \cref{lem:flag-one-dimensional-leaves} there are $\dim W_{eq}$ of them. Hence
\[
 \lim_{q\to\infty}\frac{\dim W_{eq}}{q^d}
 =\operatorname{vol}_{\mathbb R^d}\Delta'.
\]
Dividing by $e^d$ gives the asserted limit, and its positivity gives eventual nonvanishing. For $d=0$ the body is a point of volume one, so the integer dimensions converge to one and are eventually equal to one. Finally, every degree with a nonzero section lies in the projection $e\mathbb Z$ of $G$, which proves the vanishing outside these multiples.
\end{proof}

\Needspace{6\baselineskip}
The growth dimension also has a description independent of the flag. Write $\mathcal M(X)$ for the field of meromorphic functions on $X$ and define the \emph{section-ratio field}
\[
 K(W_\bullet)=\{s/t:s,t\in W_m\text{ for some }m\geq0,
                                  \ t\neq0\}\subset\mathcal M(X).
\]
It is a field: sums and products of two such ratios, with denominators $t_1\in W_{m_1}$ and $t_2\in W_{m_2}$, can be written with the common denominator $t_1t_2$ in degree $m_1+m_2$, and inversion exchanges numerator and denominator. It contains $\mathbb C$ because $W_0=\mathbb C$.

\begin{proposition}\label{prop:flag-dimension-section-field}\label{prop:full-flag-body-meromorphic}\index{Section-ratio field}
The transcendence degree of $K(W_\bullet)$ over $\mathbb C$ is exactly $d$. If $\Delta(W_\bullet)$ has nonempty interior in $\mathbb R^n$, then $X$ has $n$ algebraically independent meromorphic functions. Moreover, $d$, the period $e$ and the intrinsic volume in \cref{thm:intrinsic-flag-count} are the same for any two complete flags on $X$.
\end{proposition}
\begin{proof}
Suppose first that $d>0$. Choose sections $s_0,\ldots,s_d$ of positive degrees $m_0,\ldots,m_d$ whose normalized valuations $\nu(s_i)/m_i$ are affinely independent. Let $M$ be a common multiple of their degrees and put $t_i=s_i^{M/m_i}\in W_M$. The ratios $f_i=t_i/t_0$ belong to $K(W_\bullet)$. The flag valuation extends to quotients by subtraction, giving
\[
 \nu(f_i)=M\left(\frac{\nu(s_i)}{m_i}
                  -\frac{\nu(s_0)}{m_0}\right),
 \qquad 1\leq i\leq d.
\]
These vectors are linearly independent. Distinct monomials in the $f_i$ therefore have distinct valuations, so the term of smallest valuation in a nonzero polynomial expression cannot cancel, and the ratios are algebraically independent. Thus the transcendence degree is at least $d$; this is trivial when $d=0$.

Conversely, let $f_1,\ldots,f_r$, with $r\geq1$, be algebraically independent elements of the field. Using the product of their denominators, write $f_i=A_i/B$ with $A_i,B\in W_M$ and $B\neq0$ for a common positive degree $M$. For every integer $q\geq1$ and every multi-index $a=(a_1,\ldots,a_r)\in\mathbb N^r$ with $|a|=a_1+\cdots+a_r\leq q$, the sections
\[
 B^{q-|a|}A_1^{a_1}\cdots A_r^{a_r}\in W_{Mq}
\]
are linearly independent, because division by $B^q$ turns a linear relation among them into a polynomial relation among the $f_i$. There are $\binom{q+r}{r}$ such sections. Since $W_M\neq0$, the degree $M$ is a multiple of $e$, and \cref{thm:intrinsic-flag-count} gives $\dim W_{Mq}=O(q^d)$. Comparing with the binomial lower bound forces $r\leq d$, which proves equality of the transcendence degree with $d$. A body with nonempty interior has $d=n$, and this gives the stated consequence for $X$.

The field and the set of degrees with nonzero sections do not depend on the flag, so neither $d$ nor $e$ does. By \cref{thm:intrinsic-flag-count}, the lattice volume is, for either flag, the limit of the same dimensions with the same normalization. It is therefore independent of the flag as well, although the bodies and their lattices may differ.
\end{proof}


\section{Flags on a point blowup}
\label{sec:auxiliary-blowup-flag}

We now remove the assumption that $X$ itself carries a complete compact flag. In dimension at least two, the exceptional divisor of a point blowup carries a complete projective flag. Hartogs extension preserves holomorphic sections, so the growth formula on the modification counts the original section spaces. On a curve, a point already completes the flag.

\begin{lemma}\label{lem:point-blowup-section-descent}
Every connected compact complex manifold $X$ of dimension $n\geq1$ admits a smooth compact modification $\pi:\widetilde X\to X$ with a complete smooth compact flag. It may be chosen so that, for every holomorphic line bundle $L$ and every $m\geq0$, pullback gives
\[
 H^0(X,L^m)\simeq H^0(\widetilde X,\pi^*L^m).
\]
These identifications respect products.
\end{lemma}
\begin{proof}
If $n=1$, take $\pi$ to be the identity and choose any point as the last member of the flag. If $n\geq2$, blow up a point $p\in X$. The usual local construction of the blowup of $\mathbb C^n$ at the origin shows that $\widetilde X$ is smooth and its exceptional divisor is $E\simeq\mathbb P^{n-1}$. Choose a nested linear flag in $E$. Together with $\widetilde X$ it gives
\[
 \widetilde X\supset E\supset\mathbb P^{n-2}
       \supset\cdots\supset\mathbb P^0.
\]
Every member is smooth and compact, and is a prime divisor in the preceding member. The modification is an isomorphism off $E$.

Let $s$ be a section of $\pi^*L^m$. Off $E$, it gives a section of $L^m$ on $X\setminus\{p\}$. In a trivialization of $L^m$ on a coordinate ball at $p$, its coefficient is holomorphic on the punctured ball. The classical Hartogs extension theorem in dimension $n\geq2$ extends that coefficient holomorphically across $p$. It glues with the section away from $p$. Its pullback agrees with $s$ off $E$, and hence everywhere by the identity theorem. This proves surjectivity; injectivity and product compatibility follow from equality on the same dense open set. The identity-map case is immediate.
\end{proof}

On the modification in \cref{lem:point-blowup-section-descent}, apply the flag construction to $\widetilde W_m=\pi^*W_m\subseteq H^0(\widetilde X,\pi^*L^m)$. The resulting body's shape may depend on the point and flag; its dimension and normalized count are determined by the original section algebra.

\begin{theorem}\label{thm:all-compact-section-count}
Assume that some positive-degree piece of $W_\bullet$ is nonzero. There are an integer $0\leq d\leq n$, a degree period $e\geq1$, and a number $V(W_\bullet)>0$ such that
\[
 \lim_{q\to\infty}\frac{\dim W_{eq}}{(eq)^d}=V(W_\bullet),
 \qquad W_m=0\quad\hbox{if }e\nmid m.
\]
Here $d$ is the transcendence degree of the section-ratio field $K(W_\bullet)$, $e$ is the greatest common divisor of the nonzero degrees, and $V(W_\bullet)$ is the intrinsic lattice volume of every auxiliary body constructed above. In particular these three quantities are independent of the auxiliary modification and flag.
\end{theorem}
\begin{proof}
Pullback preserves the graded dimensions and the nonzero degrees. Apply \cref{thm:gauduchon-linear-flag-bound,thm:intrinsic-flag-count} on $\widetilde X$ to obtain the displayed limit and positive lattice volume. Pullback also identifies the section-ratio fields: their elements are the corresponding ratios, and injectivity follows on the dense isomorphism locus. Thus \cref{prop:flag-dimension-section-field} identifies $d$ with the original field's transcendence degree. The unchanged nonzero degrees determine $e$, and the limit of the unchanged dimensions determines $V(W_\bullet)$. All three are independent of the auxiliary choices.
\end{proof}

\Needspace{16\baselineskip}
\begin{corollary}\label{cor:non-moishezon-section-vanishing}
For every holomorphic line bundle $L$ on a connected compact complex manifold $X$ of dimension $n$,
\[
 h^0(X,L^m)=O(m^n).
\]
If $X$ is not Moishezon, then, for each $L$, either all its positive-degree section spaces vanish, or there is an integer $d<n$ for which
\[
 h^0(X,L^m)=O(m^d).
\]
In particular $m^{-n}h^0(X,L^m)\to0$. These conclusions apply when $X$ fails the bounded mass property.
\end{corollary}
\begin{proof}
Apply the theorem to the full section ring, taking the zero-ring-in-positive-degrees case separately. The periodic count gives $O(m^d)$ for all degrees, with $d\leq n$. If $d=n$, the section-ratio field supplies $n$ algebraically independent meromorphic functions on $X$. By Moishezon's theorem, $X$ admits a projective modification and is Moishezon. Thus a non-Moishezon manifold has $d<n$. Finally, a Moishezon manifold has the bounded mass property by \cref{lem:kahler-modification-bounded-mass}. Its contrapositive proves the last assertion.
\end{proof}

The same upper bounds apply to multiplier-ideal section spaces by their inclusion in $H^0(X,L^m)$. Identifying weighted spaces on a modification additionally requires the Jacobian correction; \cref{lem:adjoint-ideal-transform} gives that formula with the relative canonical bundle.


\section{The Hopf flag and its lower-dimensional body}

On a scalar Hopf manifold, the coordinate subspaces give a complete flag ending at a point of an elliptic fiber. Along this flag, a homogeneous section reduces to a nowhere-vanishing section on the elliptic fiber, so the last coordinate of its valuation is zero; the remaining coordinates fill a simplex whose intrinsic volume gives the section growth.

Let $X_n$ and $L_k$, with $n\geq2$ and $k\geq0$, be as in \cref{sec:hopf-homogeneous-sections}. In the covering coordinates $z_1,\ldots,z_n$, the image of $\{z_1=\cdots=z_{n-1}=0\}\setminus\{0\}$ in $X_n$ is an elliptic curve; choose a point $p$ on it. Consider the flag $Y_0=X_n$,
\[
 Y_i=\{z_1=\cdots=z_i=0\}/\langle \mathrm{e}^{-1}\rangle
 \quad(1\leq i\leq n-1),\qquad Y_n=\{p\}\subset Y_{n-1},
\]
where $/\langle\mathrm e^{-1}\rangle$ denotes the quotient of the punctured linear subspace by the contraction $z\mapsto\mathrm e^{-1}z$. Each member is smooth, compact and connected, and is a prime divisor in the preceding member. Take $W_m=H^0(X_n,L_k^m)$.

\Needspace{17\baselineskip}
\begin{proposition}\label{prop:hopf-flag-simplex}
For this flag,
\[
 \{\nu(s):0\neq s\in W_m\}
 =\{(a_1,\ldots,a_{n-1},0):a_i\in\mathbb N,
                                  \ \textstyle\sum_i a_i\leq km\}.
\]
Consequently
\[
 \Delta(W_\bullet)=
 \{(x_1,\ldots,x_{n-1},0):x_i\geq0,
                                  \ \textstyle\sum_i x_i\leq k\}.
\]
Its $n$-dimensional volume is zero. When $k>0$, its volume in the coordinate hyperplane $x_n=0$, with the Euclidean measure for which the coordinate lattice has covolume one, is
\[
 \frac{k^{n-1}}{(n-1)!}
 =\lim_{m\to\infty}\frac{\dim W_m}{m^{n-1}}.
\]
For $k=0$ the body is a point and $\dim W_m=1$.
\end{proposition}
\begin{proof}
By \cref{prop:hopf-bergman-exact}, sections of $L_k^m$ are homogeneous polynomials of degree $km$. Dividing successively by the highest powers of $z_1,\ldots,z_{n-1}$ that divide the section, and restricting at each step, leaves a nonzero multiple of $z_n^{km-a_1-\cdots-a_{n-1}}$, where $a_i$ is the $i$th valuation. This residual section has no zeros on $Y_{n-1}$, so the last valuation is zero, and the preceding ones are nonnegative with sum at most $km$.

Conversely, every such tuple is the valuation of the monomial
\[
 z_1^{a_1}\cdots z_{n-1}^{a_{n-1}}
 z_n^{km-a_1-\cdots-a_{n-1}}.
\]
The normalized values lie in the displayed simplex and, for $m=1$, include all its vertices, so their closed convex hull is the simplex. For $k>0$, its coordinate-lattice volume is $k^{n-1}/(n-1)!$, the leading coefficient of the binomial count in \cref{prop:hopf-bergman-exact}. For $k=0$, only the constant polynomials remain.
\end{proof}

\begin{figure}[htbp]
\centering
\begin{tikzpicture}[scale=1.05,>=Latex]
 \fill[blue!9] (0,0)--(3,0)--(0,3)--cycle;
 \draw[thick,blue!60!black] (0,0)--(3,0)--(0,3)--cycle;
 \draw[->] (-.25,0)--(3.65,0) node[right] {$x_1$};
 \draw[->] (0,-.25)--(0,3.55) node[above] {$x_2$};
 \node[below left] at (0,0) {$0$};
 \node[below] at (3,0) {$k$};
 \node[left] at (0,3) {$k$};
 \node[align=center] at (1,1) {area $k^2/2$};
 \node[anchor=west,align=left] at (3.5,2.5) {plane $x_3=0$\\ambient volume $0$};
\end{tikzpicture}
\caption{For the Hopf threefold and $k>0$, the body is a triangle in a coordinate plane. Its area gives the quadratic section growth, while its ambient volume vanishes.}
\label{fig:hopf-flag-triangle}
\end{figure}

